
\subsubsection{6.1 Interpretation of Key Findings}

\textbf{RLHF as a pervasive alignment signal.} The within-family LLaMA-2 natural experiment provides the cleanest available evidence: controlling for architecture and scale, RLHF fine-tuning consistently produces joint positive deltas for safety and ethics across 7B, 13B, and 70B scale points, with no counterexamples. The partial correlation $\rho_{partial}$(safety, machine\_ethics) = 0.841 (Bonferroni-significant) confirms this relationship persists after removing both confounders. The implication for evaluation practice: a practitioner who measures 6 trustworthiness dimensions independently is measuring the RLHF signal 5 times and the robustness dimension once.

\textbf{Privacy's RLHF sensitivity.} The finding that privacy is the strongest co-mover with safety ($\rho$ = 0.971) contradicts the original hypothesis, which placed privacy in the RLHF-insensitive cluster alongside robustness. Two complementary explanations are plausible: (1) RLHF reward models penalize privacy-violating outputs (PII exposure, sensitive disclosure) alongside harmful outputs --- the same annotator signal that increases safety also increases privacy; (2) TrustLLM's privacy dimension captures output-level refusal to produce PII, which mechanically overlaps with safety refusal behavior that RLHF trains. Both explanations are consistent with the data; distinguishing between them would require analysis of RLHF annotator guidelines not available for this study.

\textbf{Robustness isolation as an actionable finding.} Adversarial robustness does not co-move with the RLHF-shaped cluster. The full partial correlation ($\rho = -0.188$, p = 0.519) does not reach significance, but the scale-only control ($\rho = -0.771$, p = 0.0008) and the consistent direction of within-family deltas (all 3/3 negative) point to a real directional effect. The correct interpretation is that robustness follows a separate trajectory from the other 5 dimensions --- making it the dimension that most requires independent evaluation.

\textbf{The 3-dimension minimum evaluation set.} The MST identifies {truthfulness, fairness, privacy} as sufficient hub nodes spanning the correlation geometry. This is not a recommendation to discard other dimensions; it is a characterization of the correlation geometry: under the current TrustLLM evaluation framework, these three dimensions span most of the variance in the 6-dimensional space after confound removal.

\subsubsection{6.2 Limitations}

\textbf{Sample size (n=16).} With 2 covariates, partial correlation tests have effective df = 12 (n $-$ k\_covariates - 2), limiting power to detect moderate correlations (|$\rho$| < 0.55). The safety-robustness null finding ($\rho = -0.188$, p = 0.519) is most likely a power issue at n = 16 rather than true independence: scale-only ablation ($\rho = -0.771$, p = 0.0008) confirms the effect exists. All SUPPORTED findings (P1, P4) rely on large effects ($\rho$ = 0.841, 0.971, silhouette = 0.637) that are robust at this sample size.

\textbf{RLHF mechanism is proposed, not causally proven.} The LLaMA-2 within-family evidence provides the closest available causal design (controlling for architecture and scale), but the study is observational. We cannot rule out alternative confounders in cross-model data. The RLHF explanation is the most parsimonious given the evidence; we do not claim it is proven.

\textbf{TrustLLM-specific operationalization.} Findings are specific to TrustLLM's 2024 evaluation framework. HELM replication was planned as a scope extension but not executed in this study; cross-framework generalizability is identified as important future work.

\textbf{Machine\_ethics MST edge ambiguity.} The machine\_ethics--privacy edge appears in 68.8\% of bootstrap resamples, reflecting near-tie geometry at n=16. The minimum evaluation set \textit{size} (3 dimensions) is stable; the specific topology involving machine\_ethics is ambiguous.

\textbf{Data loading implementation deviation.} The Phase 4 experiments loaded TrustLLM scores from a hard-coded Python dictionary (values transcribed from published TrustLLM Table 1) rather than programmatically parsing TrustLLM JSON files as specified. The analysis values are identical to published TrustLLM Table 1; the deviation affects code organization, not data validity.

\textbf{Convenience sample.} The TrustLLM 16-model set is not a random sample of LLMs --- it focuses on prominent frontier models available in mid-2024. Findings strictly generalize to frontier text LLMs in the TrustLLM 2024 evaluation setting. The model set spans diverse architectures, scales (7B--175B), and alignment approaches, providing internal heterogeneity sufficient for the correlation analysis.

\subsubsection{6.3 Broader Impact}

\textbf{For benchmark design.} Future trustworthiness benchmarks should report a correlation structure summary alongside dimension scores, enabling practitioners to identify redundant dimensions within each model generation.

\textbf{For deployment evaluation.} The MST-derived minimum evaluation set {truthfulness, fairness, privacy} enables principled evaluation compression for the RLHF-shaped cluster, with adversarial robustness measured independently. Practitioners who use all 6 dimensions for comparative ranking are performing redundant measurement for 5 of 6 dimensions.

\textbf{For alignment research.} The privacy--safety coupling ($\rho$ = 0.971) suggests RLHF alignment achieves broader trustworthiness benefits than explicitly targeted --- a property with implications for constitutional AI and RLHF policy design.

\textbf{Potential negative uses.} Understanding which dimensions are correlated could theoretically assist adversarial actors in targeting models that score high on correlated dimensions while performing poorly on robustness. Making the structure explicit enables defensive countermeasures (targeted robustness evaluation) that are not possible without this knowledge.

